\section{五十川研・今井悠人さんの発表について}

\subsection{全天球画像を入力とした自由応答形式のVQAベンチマークの構築}
本発表では，全天球画像（360度画像）を用いた自由応答形式VQA（Visual Question Answering）ベンチマークの構築について報告があった．視野全体をカバーできる画像を活用することで，従来の視野制限による課題を克服する狙いがあると理解した．

\subsubsection{質問の背景と意図}
VQAタスクにおける画像の空間的取り扱いがモデル性能に与える影響に興味があり，「画像を切らない」とは具体的にどういう処理なのか確認したいと考えた．また，全天球画像を用いる利点が，複数の平面画像を用いた場合と比較してどのような点にあるのか，特に方向性・空間理解の点から明確にしたかった．

\subsubsection{質問内容}
\begin{itemize}
	\item 「画像を切らないで入力する」とは，パッチ分割を行わないという理解で正しいか？
	\item 全天球画像を必要とするシーンにおいて，複数枚の平面画像を使う構成と比べたときの利点は何か？（平面画像はDepth等との統合がしやすい印象がある）
\end{itemize}

\subsubsection{回答内容}
\begin{itemize}
	\item 「画像を切る」というのはcropを意味しており，特に極付近では歪みが大きいため単純なcropでは不十分であるとのことだった．
	\item 平面画像を複数用いる方法では，個々の画像がどの方向を表しているかを明示的に示す手段がなく，方向情報が失われる可能性がある．MLLMへの入力順序も学習に影響しうる点が指摘された．一方で，Depthなどの統合については，パノラマ形式のLiDARを扱う既存データセットもあるため，技術的には対応可能という見解が示された．
\end{itemize}

\subsubsection{考察}
360度画像の空間的な一貫性・方向性保持という強みが明確になった．今後，マルチモーダル統合を考慮する際には，単純な画像数ではなく空間構造の持続性を含めた設計が重要だと再認識した．入力の切り方（cropやpatch）と全体構造保持のトレードオフに関しても，より深掘りする価値がある．